% Motivation/method fragment for integration into the manuscript.
% Merge the accompanying bibliography keys without duplicating existing entries.
% This fragment contains no claim that the new layout-query experiment is complete.

Rule-based parasitic extraction separates process characterization from layout
 evaluation. Reference capacitances for representative interconnect geometries
are encoded in empirical models or lookup tables, and layout extraction evaluates
the corresponding models for encountered structures~\cite{tsai2025geometry,yang2023cnncap}.
OpenRCX provides a concrete example: its extraction rules are calibrated from
reference parasitics and are evaluated using coupling distance and surrounding
routing context~\cite{orcx_runtime,orcx_calibration}.

A finite set of characterization points does not coincide with every geometric
configuration encountered in a routed layout. Between sampled points, the
extractor must estimate capacitance using the available rule data. Model-library
accuracy therefore depends not only on pattern coverage but also on how the
nonlinear geometry--capacitance response is represented between characterization
points~\cite{zhou2026rmlpcap}. Increasing characterization density can improve
that representation, but also increases the amount of reference data that must
be generated and stored~\cite{yang2023cnncap}. This motivates improving rule
evaluation while holding the available characterization samples fixed.

Prior work has automated pattern and formula construction from field-solver
samples~\cite{li2020autoprechar}, developed neural capacitance
models~\cite{yang2023cnncap,zhou2026rmlpcap}, and studied geometry and
layer-thickness exploration~\cite{tsai2025geometry}. We focus on the complementary
problem of estimating uncharacterized queries from a supplied set of
process-specific characterization samples.

We formulate the task as characterization-conditioned capacitance estimation.
For a query geometry, a numerical anchor summarizes the available samples, while
a shared Transformer predicts residual corrections from geometry, material
context, and observed capacitances. Component-wise fusion combines numerical and
neural corrections, and a lightweight binary variant retains the same
conditioning interface. Learned parameters remain fixed during conditional
inference, while changes in the supplied characterization data alter the
predictions. The estimator can supply additional rule values for an otherwise
unchanged extraction flow.

We evaluate the estimator using controlled characterization subsampling,
fixed-support-budget comparisons, cross-process prediction, and native
extraction with the generated rule data. Controlled subsampling retains a set of
characterized samples and reserves the remaining reference points for evaluation,
separating scattered, contiguous, and boundary-directed coverage conditions.
Native extraction and timing analysis quantify agreement with the full-rule
reference on fixed routed designs.

% After the layout-derived off-grid experiment is actually completed, replace or
% extend the final paragraph using its measured cohort, reference semantics,
% baseline identity and results. See DAC_INTRODUCTION_AND_SOURCES.md.
